% -- ---------------------------------------------------yy
% pdflatex + report
% cf) xelatex 로 compile 해도 됨
%
% 전통적인 PDFLaTeX 환경에서 한글과 영문을 동시에 사용하면서 UTF-8 인코딩 경고나 폰트 에러를 없애기 위해 흔히 \usepackage[T1]{fontenc}, \usepackage{lmodern}, 그리고 한글 처리를 위한 \usepackage{kotex} 조합을 사용했습니다. 
% 이 방식을 사용하면 영문은 세련된 Roman(로만) 계열인 Latin Modern 폰트로 깔끔하게 출력되고, 한글 깨짐이나 폰트 누락 경고가 발생하지 않습니다.
%
% report 클래스를 기반으로 수식, 한글, UTF-8 인코딩, Latin Modern 폰트 설정이 반영된 LaTeX 템플릿
%
% TeX Live 2018 이후부터는 utf8 인코딩이 기본이므로 에러 없이 깔끔하게 컴파일됩니다.

% 컴파일러: 전통적인 pdfLaTeX으로 컴파일하시면 lmodern 패키지와 kotex가 맞물려 빠르고 에러 없이 깔끔한 PDF를 만들어 줍니다. (물론 XeLaTeX 환경에서도 작동합니다.)
%
% 장/절 표기: \usepackage[hangul]{kotex} 옵션 덕분에 "Chapter 1" 대신 "제 1 장"으로 보기 좋게 자동 변환됩니다.
% -- ---------------------------------------------------yy

\documentclass[11pt, a4paper]{report}

% 1. 인코딩 및 폰트 설정
\usepackage[T1]{fontenc}       % 유럽어 및 확장 라틴 문자 인코딩 지원
\usepackage{lmodern}          % Latin Modern 폰트 (Roman 계열 에러 방지용)
\usepackage[utf8]{inputenc}    % UTF-8 입력 인코딩 지정 (최신 LaTeX은 생략 가능하나 명시)

% 2. 한글 지원 설정
\usepackage[hangul]{kotex}     % 한글 문서 작성 패키지 ('hangul' 옵션으로 장, 절 이름 한글화)

% 3. 수식 및 수학 기호 패키지
\usepackage{amsmath}           % 다양한 수식 환경 (align 등) 지원
\usepackage{amssymb}           % 추가적인 수학 기호 및 글꼴 지원
\usepackage{amsfonts}          % 수학적 집합 기호 등 지원
\usepackage{bm}                % 수식 내 굵은 글씨(\bm) 지원

% 4. 기타 유용한 패키지 (여백 및 링크)
\usepackage[margin=1in]{geometry} % 문서 여백 설정
\usepackage{hyperref}          % 문서 내 하이퍼링크 및 PDF 목차 연동

% 문서 정보 설정
\title{pdflatex latex engine + report class :  수식 및 한글 보고서 템플릿}
\author{작성자 이름}
\date{\today}

\begin{document}

% 표지 생성
\maketitle

% 목차 생성 (선택 사항)
\tableofcontents
\newpage

% 제 1 장 시작
\chapter{서론}
2222222222222222222222222222222222
\section{연구 배경}
여기에 본문 내용을 작성하시면 됩니다. UTF-8 인코딩을 기반으로 작성되므로 한글이 깨지지 않고 깔끔하게 출력됩니다. 영문 폰트는 에러 없이 \textit{Latin Modern Roman} 글꼴로 미려하게 조판됩니다.

\section{수식 작성 예시}
수식 패키지가 포함되어 있어 인라인 수식과 블록 수식을 모두 자유롭게 작성할 수 있습니다.

본문 안의 수식(인라인 수식)은 다음과 같이 $E = mc^2$, \( a^2 \)  형태로 작성합니다. 차분 방정식이나 벡터 계산 시 자주 쓰이는 기호들도 모두 지원합니다.

따로 한 줄을 차지하는 독립 수식(블록 수식)은 아래와 같이 작성합니다:
\begin{equation}
    f(x) = \frac{1}{\sqrt{2\pi\sigma^2}} \exp\left( -\frac{(x-\mu)^2}{2\sigma^2} \right)
\end{equation}

여러 줄로 정렬된 수식을 작성할 때는 \texttt{align} 환경을 사용하면 편리합니다:
\begin{align}
    \nabla \times \mathbf{E} &= -\frac{\partial \mathbf{B}}{\partial t} \\
    \nabla \times \mathbf{H} &= \mathbf{J} + \frac{\partial \mathbf{D}}{\partial t}
\end{align}


\chapter{본론}
\section{시스템 모델}
새로운 장(Chapter)은 자동으로 다음 페이지에서 시작하며, 보고서(\texttt{report}) 양식에 맞추어 레이아웃이 정렬됩니다. 여기에 필요한 내용을 계속 추가해 나가시면 됩니다.

\end{document}
